\chapter{Intenção de compra de marcas sustentáveis da Amazônia — Uma abordagem hibrida de modelagem de equações estruturais e modelo de Rasch}


\begin{flushright}
Luiz Diego Vidal Santos  \footnote[1]{Programa de Pós-graduação em Ciências da Propriedade Intelectual-PPGPI/Universidade Federal de Sergipe-UFS. *E-mail:  \href{mailto:vidal.center@academico.ufs.br}{vidal.center@academico.ufs.br}}
\\[0.2cm] 
Francisco Sandro Rodrigues Holanda \footnote[2]{Departamento de Engenharia Agronômica/Universidade Federal de Sergipe, São Cristóvão, Brasil. E-mail \href{mailto:fholanda@academico.ufs.br}{fholanda@academico.ufs.br}}
\\[0.2cm] 
Catuxe Varjão de Santana Oliveira \footnote[3]{Departamento de Ciências da Computação/Instituto Federal de Sergipe, Lagarto, Brasil. E-mail \href{mailto:catuxe.oliveira@ifs.edu.br}{catuxe.oliveira@ifs.edu.br}}
\\[0.2cm] 
Pedro Vinícius Bertulino de Menezes \footnote[2]{Departamento de Engenharia Agronômica/Universidade Federal de Sergipe, São Cristóvão, Brasil. E-mail \href{mailto:pedrovinicius18@academico.ufs.br}{pedrovinicius18@academico.ufs.br}}
\\[0.2cm] 
Ana Paula Schervinski Villwock \footnote[2]{Departamento de Engenharia Agronômica/Universidade Federal de Sergipe, São Cristóvão, Brasil. E-mail \href{mailto:ana.agronomia@gmail.com}{ana.agronomia@gmail.com}}
\\[0.2cm] 
Alceu Pedrotti \footnote[2]{Departamento de Engenharia Agronômica/Universidade Federal de Sergipe, São Cristóvão, Brasil. E-mail \href{mailto:alceupedrotti@gmail.com}{alceupedrotti@gmail.com}}
\end{flushright}

% resumo em português
\begin{resumoumacoluna}
O artigo explora as facetas do \textit{Brand Equity} baseado em marcas relacionadas à sustentabilidade ambiental da Amazônia. A ênfase da pesquisa está nos aspectos que influenciam a predileção dos consumidores aos produtos sustentáveis oriundos de marcas em cinco categorias: cosméticos, turismo, roupas, acessórios e alimentos. Esta pesquisa planejou identificar quais os construtos do \textit{Brand Equity} baseados em marcas relacionadas à sustentabilidade ambiental da Amazônia influenciam a lealdade do consumidor à marca e intenção de compra de seus produtos. Foi conduzida uma abordagem mista de modelagem por equações estruturais e teoria de resposta ao item. Oito marcas foram selecionadas como suporte de informação, tendo como amostragens 535 consumidores. Foram considerados construtos formativos da intenção de compra PI as variáveis: posicionamento estratégico da marca — SBP, percepção da qualidade dos produtos — BQP e Inovação no varejo — RI e Autenticidade da Marca — BI. Os resultados dos modelos de equações estruturais demonstram que todos os caminhos são estatisticamente significativos e com poder de efeito forte. O posicionamento estratégico é o maior contribuinte para o efeito na intenção de compra. Como hipotetizado, o \textit{Brand Equity} sustentável ambiental tem efeitos indiretos na intenção comportamental da compra. Os consumidores apresentam um nível de $\delta$ alto, já os itens apresentam valores altos corroborando a aceitação dos produtos, exceto para diferenciação das marcas. Em geral, os resultados apoiam o uso de um \textit{Branding} voltado à sustentabilidade ambiental para balizar o \textit{Brand Equity} de pequenas empresas florestais.
\vspace{\onelineskip}
 
\noindent
\textbf{Palavras-chave}: Sustentabilidade ambiental; Marcas; Startups.

\noindent
Artigo sob avaliação no periódico: Mitigation and Adaptation Strategies for Global Change
%\href{https://doi.org/10.53660/CLM-106-127B}{\color{blue}{\underline{https://doi.org/10.53660/CLM-106-127B}}}

\noindent

\begin{figure}[H]
  \centering
  \caption*{Resumo gráfico 3º capítulo}
  \includegraphics[scale=0.15]{Imagens/resumo_3.pdf}
  \fonte{\textbf{Elaborado pelos autores} (2022)}
  \label{figres:04}
\end{figure}
\href{https://empreendasustentavel.com/resumo_3.html}{\color{blue}{\underline{Download: Baixe a imagem em tamanho real}}}


\end{resumoumacoluna}
\section{Introdução}

O agravamento dos problemas ambientais está modificando a maneira como o conhecimento técnico-científico aplica-se no processo produtivo. Os danos ao meio ambiente não são fatos inesperados e imprevisíveis, eles apenas demonstram a urgência do conhecimento científico minimizar os efeitos gerados pelo desenvolvimento industrial \cite{Adger2009NestedChange,Chen2021InvestigationGrowth}. Nesta perspectiva, as ações para conservação ambiental são vistas como importantes fatores da contribuição para o desenvolvimento sustentável de um país \cite{Lamichhane2021BenchmarkingApproach}. 

Se de um lado a revolução industrial apresentou um novo cenário para as atividades econômicas e humanas, os novos processos e as novas tecnologias transformaram as concepções sobre os recursos naturais, o ser humano e o planeta para além de paradigmas estáticos \cite{Cabigiosu2022SustainableVehicles}. A natureza multidisciplinar do desenvolvimento sustentável e as dificuldades em reconhecer seu significado dificultam muito seu monitoramento por meio de indicadores únicos. Neste sentido, é necessário criar estruturas integradas para considerar plenamente seus aspectos sociais, econômicos e ambientais interdependentes. A ocorrência se dá pela dificuldade em medir elementos isolados do desenvolvimento sustentável com quaisquer indicadores únicos \cite{Diaz-Sarachaga2018IsAgenda}. Visando compreender e desenvolver estratégias além dos indicadores estáticos e únicos, a Assembleia Geral da ONU em 2015 criou os Objetivos do Desenvolvimento Sustentável - ODS (Agenda 2030).

A agenda 2030 foi criada como resultado de anos de discussões envolvendo legisladores, acadêmicos, funcionários públicos (inter) nacionais, elites políticas e grupos de base, representando 193 países membros da ONU \cite{ONU2015}. Os Objetivos de Desenvolvimento Sustentável (ODS) das Nações Unidas enfatizam o importante papel de vários atores (não estatais), incluindo empresas e ONGs, para garantir a sustentabilidade da cadeia de abastecimento e reduzir o desmatamento tropical e a degradação florestal \cite{Delabre2020StrategiesGoals}.

Os 17 objetivos elencados foram desenvolvidos por meio de processos consultivos massivos que incluíram todos os envolvidos nos aspectos sociopolíticos, econômicos e ambientais do desenvolvimento \cite{Novovic2022CanSystem}. Mecanismos para atingir os objetivos do desenvolvimento sustentável se baseiam em abordagens integradas e multissetoriais, incluindo uma abrangente participação de diversos atores, públicos ou privados, promovendo-se uma sinergia entre dissemelhantes frentes de meios de produção e consumo: agricultura, alimentação, saúde, nutrição, educação, desenvolvimento social, economia, entre outros \cite{Silva2018g}. 

A interação ciência, tecnologia e sustentabilidade requer uma sinergia sustentada por paradigmas que incorporam o potencial ecológico, as condições ambientais e os valores culturais na organização dos processos produtivos \cite{Garcia-Muina2018TheCompany,Morteza2019Eco-capabilityParadigm}. Essa sinergia requer uma promoção sustentável tanto na manutenção da economia quanto na sustentabilidade ambiental. Para isso, é preciso estar sempre em pauta o combate às práticas de degradação ambiental. Este enfrentamento deve envolver as formas de produção agropecuária, o extrativismo florestal e os seus sistemas, tornando tais temas pautas frequentes nas empresas e em suas políticas de governança.

Buscando mitigar os desafios ambientais, pesquisadores apresentam novas possibilidades de lidar com a complexa rede atual de interdependência do desenvolvimento e a sustentabilidade por meio da inovação científica e tecnológica \cite{Zeigermann2020ChallengesDevelopment}. Sob essa perspectiva, pesquisas tem se pautado no empreendedorismo ambientalmente sustentável para impactar o desenvolvimento para além do âmbito econômico, mas com impactos sociais e ambientais, por meio do processo de descobrir novas oportunidades de negócios inovadores e sustentáveis com os seus enfoques distintos \cite{Sepulveda2011, Davidsson2016}.

O empreendedorismo com foco na sustentabilidade ambiental pode maximizar lucros de oportunidades de negócios, contribuindo para redução da pobreza e impactando positivamente na sociedade e no meio ambiente \cite{Iqbal2021InsightsPerformance}. No entanto, para que o empreendedorismo ecologicamente responsável seja implementado pelas empresas de forma satisfatória, é essencial que seja economicamente viável, o que promoverá uma maior aceitação das empresas orientadas para esse propósito. Nesse aspecto, o gerenciamento de marcas com foco no comportamento sustentável da política empresarial se mostra necessário.

Estratégias de gerenciamento de marketing e negócios concebidos nos princípios do desenvolvimento sustentável serão importantes na definição de tecnologias-chave onde as empresas devam investir tanto para a resolução dos seus problemas ambientais básicos quanto para uma política de exportação de produtos, principalmente para as empresas que prezam a sustentabilidade ambiental \cite{Lewis2012107, Tugrul2017}. Atualmente, há um consenso de que a aplicação de capital financeiro em pesquisa é um investimento fundamental para o desenvolvimento sustentável e para a melhoria da qualidade de vida das populações \cite{Nevado-Pena2019ImprovingEurope,Ikram2021AssessingContext}. As pesquisas em gerenciamento de marketing e empreendedorismo devem apresentar soluções que tenham como objetivo controlar os problemas socioambientais provocados pelo desenvolvimento predatório, ao mesmo tempo que mantêm o empreendedorismo economicamente competitivo \cite{Loh2020,Liu2020b}. Uma das áreas de pesquisa no gerenciamento de marketing é a percepção da experiência do consumidor em relação à marca, também conhecida como experiência da marca.

Experiências da marca podem ser classificadas como subjetivas, internas (sensações, sentimentos e cognições) e respostas comportamentais evocadas por estímulos relacionados à marca, e compõem o design, identidade, embalagem, comunicação e ambientes de uma marca \cite{Ramaseshan2014ConnectingRelationships}. Os Estímulos oriundos das experiências da marca podem ser quaisquer pistas experienciais que desencadeiam dimensões de experiência da marca, levando a melhores relações de marca em qualquer ponto de interação \cite{Jeon2017TheEquity}. 

O desenvolvimento sustentável e a responsabilidade social assumem um papel cada vez mais importante na experiência da marca \cite{Ketprapakorn2019}. Buscando associar tais fatores, estudos analisaram a intenção de compra da marca a partir de perspectivas culturais, sociais, psicológicas e econômicas \cite{Cheah2015AntecedentsBrands,Kautish2020InfluenceConsumption,Dhaliwal2020TheAgenda}, sendo estes construtos baseados nos fatores: confiança, satisfação, lealdade e compromisso com a marca \cite{Kim2012DoBrand,Zollo2020UnpackingExperience}. 

A categorização de marca sustentável apresenta fatores além dos construtos citados, tais como o apego emocional, a consciência social e ambiental, assim como o senso de exclusividade nos produtos \cite{Long2018366,Chen20191,Ranfagni2022LuxuryOpportunities}. O consumo de produtos ecologicamente responsáveis envolve uma forte relação atitudinal e de apoio mútuo entre o consumidor e a marca \cite{Cheng2022TheAttachment}, muitas vezes ocorrendo por antropomorfização da marca, dotando-a de sentimentos e responsabilidades comuns \cite{Langner2015IsLove,Schmid2019BrandLifecycle}.

Pensando nisso, várias estratégias de marketing recorrem aos estímulos psicológicos para propagar a imagem da marca no mercado consumidor \cite{Yoo2017AnEquity:}. Algumas utilizam táticas de singularidade nos produtos, outras utilizam a qualidade percebida nos produtos sustentáveis como diferencial \cite{Liu2014ThePreference} ou a disposição pessoal de adquirir o produto por valores éticos agregados \cite{Pero2020SustainabilityBrands}, sem contar com a exposição da inovação ecologicamente sustentável utilizada no desenvolvimento dos produtos \cite{Nedergaard2013SustainableInnovation}.

Considerando a importância de estudos desta natureza, para a definição de prioridades e fomento no campo dos investimentos em marketing, políticas governamentais relacionadas ao desenvolvimento empresarial sustentável e tecnológico do país e implementação e avaliação no nível local de viabilidade negocial, esta pesquisa objetiva identificar quais os construtos do \textit{brand equity} baseados em marcas relacionadas à sustentabilidade ambiental da Amazônia influenciam a intenção de compra dos consumidores ao produto.

\section{Referencial teórico}



\section{Material e métodos}

Foram investigadas oito marcas relacionadas que comercializam produtos que utilizam a marca Amazônia sustentável. As marcas são de cinco categorias de produtos diferentes: cosméticos, turismo, roupas, acessórios e alimentos. Para caracterização da percepção da marca, reunimos produtos, embalagens, \textit{mockups} de produtos disponíveis nos sites das marcas, \textit{posts} em redes sociais e \textit{copy} para Marketing Digital.

Foram entrevistados 535 potenciais consumidores para esta pesquisa. Esse tamanho de amostra atende ao limite mínimo recomendado por \citeonline{Hair2017AnResearch}, segundo o qual o tamanho da amostra deve ser igual ou maior que: (i) dez vezes o número de indicadores formativos usados para avaliar um construto, ou (ii) dez vezes o maior número de caminhos estruturais direcionados a um construto latente específico no modelo estrutural \cite{Hair2017AnResearch}. 

Devido à ausência de um quadro de amostragem de todos os consumidores das marcas, bem como para evitar enviesamento da percepção, na seleção dos entrevistados participantes, foi utilizada a amostragem não probabilística. Para esta pesquisa, foram aplicados questionários on-line e presenciais durante os meses de março a setembro de 2022. Após os respondentes preencherem os formulários da pesquisa, eles foram recolhidos \textit{in loco} pelos pesquisadores, transcritos digitalmente, preservando potenciais dados.



\subsection{Construtos analisados}

O constructo intenção de compra (\textit{Purchase intention} - PI) contém 5 itens que buscam compreender o quanto os participantes estariam dispostos a adquirir produtos de empresas relacionadas à sustentabilidade ambiental por meio dos dados apresentados. Para medir a intenção para compra, foram adaptados 5 (cinco) itens dispostos no trabalho de  \citeonline{Arnett-meiers}.

O segundo constructo foi inicialmente composto por questões transliteradas da escala desenvolvida por \citeonline{Grigorescu2019KeySustainability}, que avalia a percepção do entrevistado sobre o posicionamento estratégico de marcas sustentáveis. Dos 21 (vinte e um) itens da escala original, foram utilizados nesta pesquisa 6 (seis) com pequenas modificações de linguagem. Adotando, a partir das modificações, a identificação para o constructo de Posicionamento Estratégico de Marca \textit{(Strategic Brand Positioning — SBP)}.

O terceiro constructo contou com 5 itens sobre a Percepção da qualidade da marca (\textit{Brand Quality perception} — BQP) e o \textit{Brand Equity} relacionado à sustentabilidade ambiental do instrumento proposto por \citeonline{Netemeyer2004DevelopingEquity}. Este construto foi avaliado utilizando cinco figuras e itens da marca, tanto on-line quanto presencialmente, cinco imagens da loja (digital), quatro figuras ou itens do produto \textit{premium} (com valor acima da média do mercado) e cinco figuras que representaram a intenção empresarial frente à sustentabilidade ambiental da Amazônia (meio de produção/extração)

O quarto constructo avaliou o quanto o negócio consegue apresentar sua inovação no varejo (\textit{Retail Innovation} — RI) e ser reconhecido por meio do \textit{Branding}. A dimensão foi montada solicitando aos participantes que respondessem a 3 questões que abordam o tema inovação no varejo. As questões foram adaptadas dos questionários propostos por \citeonline{Lin2015ConceptualizingTaiwan} e \citeonline{Netemeyer_Krishnan}.

O quinto constructo avaliou o quanto o negócio consegue ser autêntico para os consumidores (\textit{Brand Authenticity} — BA) e ser reconhecido por meio do \textit{Branding}. O BA foi medido com duas perguntas baseadas nos trabalhos de \citeonline{Lin2015ConceptualizingTaiwan} e \citeonline{Spiggle2012MoreAuthenticity}.
\subsection{Instrumento utilizado}

Para obtenção dos dados da pesquisa, foi utilizado o instrumento BES-Amazônia, consumidores traduzidos e revalidados (Apêndice H), tendo como base para validação modelo tripartite, descrito no \textit{Standards for Educational and Psychological Testing}\cite{Association2018}, que apresenta o processo de validação nas seguintes fases: 

\begin{enumerate}[label=(\roman*)]
- Validação de conteúdo;
- Validade baseada na estrutura interna;
- Validade baseada nas relações com medidas externas;
- Validade baseada no padrão de resposta aos itens.
\end{enumerate}

Após a validação da escala, o instrumento final foi composto por 22 itens a tratar sobre o \textit{Brand Equity} de marcas relacionadas à sustentabilidade ambiental da Amazônia. 

\subsection{Modelo estrutural hipotético}

Com base em estudos de revisão, foi hipotetizada uma rede nomológica apresentada na Figura \ref{fig:12}. O modelo propõe compreender a disposição do consumidor em adquirir o produto e sua lealdade para com a marca. Sendo assim, foram analisados quatro construtos (i) intenção de compra PI (5 questões), (ii) Percepção da Qualidade — BQP (5 questões); (iii) Percepção do posicionamento estratégico da marca — SBP (8 questões) a (iv) Percepção da inovação no varejo — RI (3 questões) e a (v) Autenticidade da marca BA (2 questões). Foram testadas quatro hipóteses de pesquisa tendo como base os construtos estudados:


- ($H_1$) A percepção da qualidade tem um efeito positivo na intenção de compra de marcas sustentáveis;
- ($H_2$) O posicionamento estratégico da marca tem um efeito positivo na intenção de compra de marcas sustentáveis;
- ($H_3$) A inovação da marca no varejo tem um efeito positivo na intenção de compra de marcas sustentáveis;
- ($H_4$) A autenticidade na marca tem um efeito positivo na intenção de compra de marcas sustentáveis.


\begin{figure}[H]
  \centering
  \caption{Rede nomológica de relacionamentos conceituais da percepção estudada}
  \includegraphics[scale=0.5]{Imagens/fig_12.pdf}
  \fonte{\textbf{Elaborado pelos autores} (2022)}
  \label{fig:12}
\end{figure}

\subsection{Análise dos dados}

Os dados da escala foram inicialmente avaliados quanto à normalidade usando o teste de Kolmogorov-Smirnov. Como os dados não apresentavam distribuição normal, foi realizada a correlação de Spearman. Não foi encontrada multicolinearidade.

Um modelo de segunda ordem foi proposto utilizando a Análise Fatorial Confirmatória AFC. Inicialmente, para a implementação da Modelagem por Equações estruturais (MEE), foi definido um modelo que evidencia as relações hipotéticas entre o conjunto de variáveis analisadas decorrentes da obtenção de relações empíricas resultantes da AFE. Para representar tais relações, é frequentemente utilizado um esquema visual denominado Diagrama de Caminhos. 

Para analisar ainda mais os resultados correlacionais, uma modelagem de equações estruturais foi implementada utilizando o software R, por meio do método dos \textit{Weighted Least Squares Mean and Variance Adjusted} (WLSMV) similar ao RDWLS, útil para dados categóricos presentes em matrizes policóricas não normalmente distribuídos. Os efeitos de mediação foram obtidos pela técnica de \textit{bootstrapping} com 5.000 reamostragens para um intervalo de confiança de 95\% do efeito mediado. 
Para avaliar o modelo global, foram considerados os mesmos índices de ajuste utilizados na AFE (RMSEA, CFI e TLI), assim como o $\chi$2 e $\chi$2/gl. Valores de $\chi$2 não devem ser significativos; a razão $\chi$2/gl deve ser < que 5 ou, preferencialmente, < que 3 \cite{Kline2016PrinciplesLivros}. Valores de CFI e TLI devem ser > que 0.90 e, preferencialmente, acima de 0.95 \cite{Sellbom2019FactorRecommendations}. Valores de RMSEA devem ser < que 0.08 ou, preferencialmente, < que 0.06, com intervalo de confiança (limite superior) < 0.10 \cite{Brown2015}. A fidedignidade da medida foi mensurada através da fidedignidade composta \cite{Raykov2016EstimationMeasures}.

A acurácia preditiva dos construtos foi avaliada pelo coeficiente de determinação ($R_2$). Segundo \cite{HairJr2021APLS-SEM}, nas pesquisas que se baseiam no comportamento do consumidor, valores de $R_2$ de 0.75, 0.50 e 0.25 podem ser considerados substancial, moderado e fraco, respectivamente.

A análise dos itens foi proposta utilizando o modelo de teoria de resposta ao item com dois parâmetros logísticos; método Rasch. As dificuldades dos itens do modelo foram avaliadas por meio do modelo Rasch para dados politômicos \cite{Andrich1978ACategories}. Foram avaliados indicadores da fidedignidade das pessoas e dos itens, bem como desvios de desempenho, por meio dos índices de \textit{Infit} e \textit{Outfit}. O DIF foi avaliado por meio do procedimento de Mantel \cite{Bond2013ApplyingSciences}. Itens cujas estimativas de dificuldades se apresentaram como estatisticamente diferentes para homens e para mulheres (p$\leq$ 0.05) foram inspecionados. As análises foram desenvolvidas utilizando o software Winstep versão 5.1.4.

\section{Resultados e discussão}
\subsection{Estatísticas descritivas}

A demografia dos entrevistados inclui 52\% do sexo masculino e 48\% do sexo feminino de todas as regiões do Brasil. Além disso, 25\% dos entrevistados tinham idade entre 15 e 24 anos, 23\% entre 25 e 34 anos, 24\% entre 35 e 44 anos, 22\% entre 45 e 54 anos e 6\% entre 55 e 64 anos. Quanto à renda mensal, 40\% renda mensal domiciliar até 555 US\$ 14\% possui renda mensal domiciliar entre 555 US\$ e 1358 US\$ e 12\% possuem renda mensal domiciliar entre 1358 US\$ e 4210 US\$.

\subsection{Estrutura interna e resultados de confiabilidade para o modelo.}

A estrutura unidimensional proposta se ajustou bem aos dados. Embora o valor do c2 tenha sido significativo, todos os outros índices de ajuste suportaram o modelo (Tabela \ref{tab:13}), corroborando assim o ajuste do modelo aos dados \cite{Steiger2010StructuralApproach,Bagozzi2016StateUsage}

\begin{table}[H]
\centering
\caption{Índices de ajuste do modelo bifatorial}
\label{tab:13}
\begin{tabular}{llllll}
\hline
\multicolumn{1}{c}{$\chi_2$ (gl)} &
  \multicolumn{1}{c}{$\chi_2$/gl} &
  \multicolumn{1}{c}{CFI} &
  \multicolumn{1}{c}{TLI} &
  \multicolumn{1}{c}{SRMR} &
  \multicolumn{1}{c}{RMSEA (90\% IC)} \\ \hline
904.603(183) &
  4.94 &
  0. 910 &
  0.981 &
  0.781 &
  0.084 (0.071 — 0.109) \\ \hline
\end{tabular}

\footnotesize{Nota: $\chi_2$ = qui-quadrado; gl = graus de liberdade; CFI = Comparative Fit Index; TLI = Tucker-Lewis Index; SRMR = Standardized Root Mean Square Residual; RMSEA = Root Mean Square Error of Approximation.}
\end{table}

As medidas de confiabilidade e validade são apresentadas na Tabela \ref{tab:14}. Todos os itens apresentaram valores adequados, com cargas fatoriais elevadas em seus respectivos construtos.

\begin{table}[h]
\centering
\caption{Medidas de confiabilidade do instrumento}
\label{tab:14}
\begin{tabular}{lllll}
\hline
\multicolumn{1}{c}{}  & \multicolumn{1}{c}{SBP} & \multicolumn{1}{c}{BQP} & \multicolumn{1}{c}{RI} & \multicolumn{1}{c}{PI} \\ \hline
Composite Reliability & 0.937                   & 0.703                   & 0.714                  & 0.760                   \\
H- latent             & 0.965                   & 0. 916                  & 0. 916                 & 0. 905                  \\ \hline
\end{tabular}
\end{table}

Os pesos ($\beta$) para cada um dos itens em seu fator pretendido foram significativos e todos os coeficientes padronizados > 0.40, com exceção ao item PI 03 que apresentou dados marginalmente significativos; sugerindo que cada um dos itens permaneça no modelo (Figura \ref{fig:13}).

\subsection{Testando a teoria proposta para o modelo de intenção de compra do consumidor.}

Com base nos resultados da simulação em Modelagem por Equações Estruturais (MEE) e no conjunto de variáveis observáveis, foi estabelecido o modelo de definição do \textit{Brand Equity} baseado em marcas relacionadas à sustentabilidade ambiental da Amazônia que influenciam a intenção de compra dos consumidores ao produto.

O modelo estrutural (Figura \ref{fig:13}) demonstrou que a predileção dos consumidores aos produtos de marcas sustentáveis é influenciada direta e positivamente pelas variáveis posicionamento estratégico da marca POM: $\beta$ = 0,901; erro padrão (SE) = 0.028; (p = 0.00). A pesquisa prevê, em $H_1$, que a percepção da qualidade tem um efeito positivo na intenção de compra de marcas sustentáveis, neste sentido, os resultados do estudo suportam esta proposição assumindo que, o valor de $\beta$ é satisfatório para acatar a hipótese.

Para a hipótese $H_2$ definida pelo construto percepção da qualidade dos produtos PQM: $\beta$ = 0.751; SE = 0,048; \textit{p} < 0.000, os resultados mostraram um efeito indireto significativo, assim como os caminhos padronizados entre o fator e os indicadores estão todos acima de 0.40, portanto sustentável \cite{Wang2019StructuralModels}. Por fim, a hipótese $H_3$ Percepção da inovação (RI) também apresenta efeito indireto significativo para a percepção do consumidor RI: $\beta$ = 0.691; SE = 0,070; p < 0.000.

 \begin{figure}[H]
  \centering
  \caption{Contribuição das variáveis posicionamento estratégico da marca (SBP), percepção da qualidade dos produtos (BQP) e Inovação no varejo (RI) na intenção de compra dos produtos oriundos de marcas sustentáveis da Amazônia (PI). ($\chi_2$/gl=4.94; CFI=0.910; TLI=0.981; SRMR =0.781; RMSEA=0.084 (0.071 — 0.109)}
  \includegraphics[scale=0.3]{Imagens/fig_13.pdf}
  \fonte{\textbf{Elaborado pelos autores} (2022)}
  \label{fig:13}
  \footnotesize{Nota: A figura é um gráfico representacional. Os itens que compõem cada dimensão não são mostrados para manter o gráfico simples e parcimonioso.}
\end{figure}

As análises iniciais demonstram índices adequados para os itens (Confiabilidade = 1.00; Índice de Separação = 18.77) e para as pessoas (Confiabilidade = 0.88; Índice de Separação = 2.70), o que sugere que as estimativas obtidas tendem a ser replicáveis. Em relação aos desvios de desempenho (\textit{Infit} e \textit{Outfit}), é possível perceber que o instrumento apresenta valores aceitáveis \cite{Bond2013ApplyingSciences,Boone2016RaschHow} tanto para os itens (Infit MNSQ =0.99, ZSTD = -0.50; {Outfit}) MNSQ = 1.06; ZSTD = -0.00) quanto para os participantes (\textit{infit} MNSQ =1.03, ZSTD = -0.10; Outfit MNSQ = 1.06; ZSTD = -0.10). Ao analisar as respostas de um instrumento a diferentes itens, é possível observar que os dados suportam a precisão e a confiabilidade dos dados coletados por meio da TRI. \citeonline{Sireci2021UsingTest} afirma que os dados do TRI produziram resultados confiáveis que apoiam a validade e a confiabilidade da teoria. Os resultados também destacam que as medidas de Confiabilidade \textit{Infits} e \textit{Outfits} confirmam a confiabilidade do instrumento. Não foram identificados efeitos de piso ou teto nas respostas fornecidas pelos participantes.

O mapa pessoa-item (Figura \ref{fig:14}) demonstra que os consumidores apresentam, em geral, um nível de traço latente maior do que a dificuldade dos itens (média de $\theta$ = 0.67; média da dificuldade dos itens = 0). Conforme pode ser visto, os itens que requerem maior nível de traço latente ($\theta$) foram BA 02 \textit{“Após ver a imagem acima, você conseguiria diferenciar os serviços dessa empresa de outras empresas?”} ($\theta$= 3.58) e RI 03 \textit{“Você acredita que esse símbolo/logotipo da marca é fácil de ser lembrado?”} ($\theta$= 3.40). Os itens com maior dificuldade (BA 02 e RI 03) demostram que os potenciais consumidores precisam de um traço latente maior para confirmação positiva, esse resultado mostra que, para os consumidores, os sinais visuais utilizados pelas marcas não refletem o posicionamento estratégico buscado, assim como apresentam certa dificuldade de serem lembrados posteriormente.

Analisando as características dos participantes (Figura \ref{fig:14}), é possível observar que a maioria apresenta um nível de $\beta$ alto agrupados entre 0.0 e 3.0. Os dados demonstram a aceitação e aquisição de produtos, assim como a percepção do consumidor, considerando que a maioria dos potenciais consumidores participantes compreenderam a apresentação do produto e estariam propensos a adquirir os produtos apresentados conforme as dimensões estudadas. A percepção ambiental se refere a todo o processo mental de interação de um indivíduo com seu entorno. Isso é realizado mediante vários métodos perceptivos, como a criação de valores de paisagem para cada indivíduo. Neste sentido, as marcas a partir da compreensão dos fatores sustentáveis podem aproveitar-se deste fator para agregar valor aos seus produtos, considerando o alinhamento adequado do discurso voltado a sustentabilidade \cite{Chuvieco2018FactorsUAE}.

A intenção de compra de produtos sustentáveis da Amazônia também se sustenta quando é observado o agrupamento dos participantes com possibilidade de endosso ($\beta$) positivo nestas questões, em que cerca de 98\% dos participantes apresentam 50\% de chance de concordar ou concordar plenamente.

\begin{figure}[H]
  \centering
  \caption{Mapa pessoa-item}
  \includegraphics[scale=0.8]{Imagens/fig_14.pdf}
  \fonte{\textbf{Elaborado pelos autores} (2022)}
  \label{fig:14}
\end{figure}

A Figura \ref{fig:14} demostra também que boa parte dos itens estão agrupados entre – 1.0 e aproximadamente 0.8 \textit{logit}, neste sentido é possível afirmar que o instrumento utilizado apresenta uma maior precisão na estimativa do \textit{theta} dos participantes alocados nesta amplitude do continuum e que não há comportamento discrepante ou de efeito de piso e teto.

\subsection{Curvas Características dos Itens}

As curvas características do item (ICCs) e as funções de informação do item (IIFs, linhas tracejadas) são apresentadas na Figura \ref{fig:15}. Essas curvas mostram que a informação do item não é uma quantidade estática, mas condicional aos níveis de dificuldade do item. Do ponto de vista da opinião dos consumidores, isso se traduziria em: quanto menor a dificuldade do item, maior aceitação para com os produtos. Isso se reflete em uma maior chance de os consumidores responderem concordo e concordo plenamente da escala.

As questões que apresentaram menor traço latente ($\theta$) para o endosso foram o interesse na aquisição: PI 01 — O consumidor compraria essa marca? Qual a predisposição à indicação dos produtos? PI 02 — O consumidor indicaria esta marca a um conhecido e na interpretação e compreensão da estratégia de \textit{Branding} utilizadas pelas marcas: SBP 02 — \textit{"A marca é amigável"} e SBP 03 — \textit{"A marca transmite confiança sobre sustentabilidade ambiental"}

\begin{figure}[H]
  \centering
  \caption{(a) Curvas características dos Itens e (b) Funções de informação dos itens}
  \includegraphics[scale=0.3]{Imagens/fig_15.pdf}
  \fonte{\textbf{Elaborado pelos autores} (2022)}
  \label{fig:15}
\end{figure}

Algumas considerações também podem ser feitas em relação aos ICCs: (1) Para a maioria dos itens, a probabilidade a 50\% de responderem a uma categoria de resposta satisfatória (acima de 4) está posicionado 2 logit abaixo da média $\theta$, corroborando com a aceitação positiva para quase todos os itens apresentados no questionário.

(2) O item BA 02 apresenta maior avanço para o nível de endosso 2.5 logit acima da média $\theta$ dos itens sobre as respostas positivas ao questionário (concordo e concordo totalmente), demostrando que, após ver a imagem acima, os consumidores sentem dificuldade de diferenciar os serviços da empresa analisada de outras empresas concorrentes, e o item RI 03, que representa a facilidade ou dificuldade de lembrar o símbolo/logotipo da marca testada. 

Estes dados indicam que os signos distintivos das empresas analisadas não são suficientemente informativos para definir e posicioná-las na memória do consumidor exposto brevemente, ou seja, não foi proporcionado um tempo maior para os respondentes terem mais capacidade de perceber os traços de responsabilidade ambiental das instituições analisadas, refletindo uma necessidade de melhoria para com a apresentação da entidade, seja ela nominal ou simbólica. Essa descoberta pode ser sustentada pela teoria da interferência de associações de marca descrita por \citeonline{Meyers-Levy1989TheMemory}. Esta teoria postula que a memória para as informações de uma corporação é aprimorada por associações relacionadas ao nome da marca (porque cada associação representa uma possível pista de recuperação) e que nomes compostos de palavras encontradas com frequência na linguagem provocam processamento não distintivo, assim a apresentação de signos atrelados a nomes de alta frequência pode interferir na memorização do cliente para com a marca. \citeonline{Robertson1989StrategicallyCharacteristics} afirma ainda que os nomes das marcas devem ser simples, distintivos, significativos, emocionais, recorrer a morfemas, fonemas, aliterações, consonâncias e devem fazer um som associado à classe do produto, além de serem legalmente protegidos.

É importante ressaltar que os resultados obtidos pelas questões e BA 01 ($\theta$ = 0.40) \textit{“Para mim a marca é única quando comparada com outras marcas”} e BA 02 ($\theta$ = 3.58) \textit{“Após ver a imagem acima, você conseguiria diferenciar os serviços dessa empresa de outras empresas?”}, também lançam luz sobre a capacidade de autenticidade e diferenciação das marcas estudadas.

Mesmo os participantes demonstrando que para eles a marca é única quando comparada com outras marcas, eles não conseguem diferenciá-las. Neste sentido, a apresentação da inovação contida nas marcas, mesmo que autêntica, pouca influência na lembrança e associações da marca dos consumidores pelas marcas sustentáveis da floresta amazônica, visto que 38\% dos consumidores afirmaram não saber opinar quanto à inovação contida no produto. 

Ao que se observa, as marcas analisadas não apresentaram boa capacidade de diferenciação a partir dos signos e embalagens utilizadas, já que estas se limitam aos aspectos funcionais e dados expressos em embalagens e dados informacionais dos produtos verdes que, nos signos apresentados, dificultam a formação do conhecimento da marca (Brand Awareness), importante dimensão da marca \cite{Keller2019StrategicProcess}. \citeonline{Granato2022TheSustainability} observaram que, ao considerar a embalagem como um fator ambiental relevante, os consumidores também consideram a marca e o produto como superiores em termos de benefícios, uma vez que a presença de um material biodegradável ou reciclável expresso em signos e embalagens demonstra positivamente outros benefícios, tais como a conveniência ou atratividade da sustentabilidade ambiental para os consumidores.

Em resumo, utilizar isoladamente o discurso de inovação na estratégia de \textit{Branding} pode não surtir efeito de atração de consumidores. Neste sentido, os resultados sugerem que, para as marcas analisadas alcançarem elos mais fortes na memória dos consumidores, se faz necessária uma adequação nominal para nomes que apresentem uma baixa frequência, uma maior persuasão para com os clientes quanto ao discurso contido nas impressões de embalagens e demais materiais midiáticos \cite{Siegrist2015FactorsStudy,Vincenzi2018AssessmentBrazil} e o desenvolvimento de embalagens sustentáveis, a exemplo as embalagens que utilizam artifícios de experiência continuada durante e após o consumo dos produtos. 

\subsection{Compreensão do \textit{Brand Equity} sustentabilidade da Amazônia por questões}

Os resultados sobre a percepção do que torna a marca sustentável conforme os consumidores pesquisados são apresentados na Figura \ref{fig:16}. Verifica-se que 36\% dos entrevistados concordam e 33\% concordam totalmente que a marca apresentada é mais responsável ambientalmente que a concorrência (SBP 06), compreendendo também que a marca por meio de seus produtos, signos e embalagens aparentemente cumpre o que promete (SBP 04), aproximadamente 44\% concordam e 24\% concordam totalmente.

Quanto aos valores acrescentados aplicados nos produtos, apenas 16\% concordam em pagar o valor máximo de acréscimo (20\% a mais no preço), tendo a maioria dos consumidores no estrato entre 10\% e 15\% de valor \textit{premium}. Este acrescimento se sustenta quando observamos a questão PI 04, em que o consumidor é questionado sobre estar disposto a pagar mais por produtos da marca do que por produtos vendidos por outras marcas da categoria. 32\% destes consumidores concordam com esta afirmativa. Os resultados obtidos corroboram os achados por \citeonline{Dwivedi2018BrandUniqueness} que, ao analisarem a credibilidade das marcas como fator mediador da disposição da compra com acréscimo de preço, verificou ser significativamente relevante a correlação entre a percepção positiva da marca e a disposição dos consumidores em pagar um preço \textit{premium}.

\begin{figure}[H]
  \centering
  \caption{Resumo da distribuição de pontuações dos consumidores quanto à percepção da sustentabilidade por questões respondidas}
  \includegraphics[scale=0.1]{Imagens/fig_16.pdf}
  \fonte{\textbf{Elaborado pelos autores} (2022)}
  \label{fig:16}
\end{figure}

No que tange à apresentação dos canais digitais de comunicação das marcas, os consumidores em sua maioria concordam totalmente que os canais de atendimento tais como sites e redes sociais apresentam um design criativo (RI 01 = 46\% Concordam totalmente). Tais resultados apontam que o design dos canais de atendimento pode influenciar positivamente para o \textit{Brand Equity} de marcas estudadas. \citeonline{Peng2017TheShopping} sugerem que os consumidores que consideram um produto visualmente atraente também percebem o site do \textit{e-commerce} como útil, resultando em prazer na experiência de compra e maior probabilidade de formar atitudes positivas em relação à marca. A percepção positiva do ambiente \textit{e-commerce} se mostra importante para as marcas estudadas, pois ao se perguntar o quanto os consumidores têm de facilidade para imaginar uma loja física, 37\% responderam que não saber opinar, demonstrando que o site afeta as variáveis orgânicas por meio do efeito mediador do apelo visual percebido, influenciando mais positivamente que o meio físico, resultado também observado por \cite{Liu2013WebsitePerceptions}

\section{Conclusões}

Esta pesquisa fez progressos significativos na compreensão da relação entre consumidores e marcas sustentáveis, especialmente no contexto da Floresta Amazônica. O estudo identificou construtos-chave que influenciam a lealdade do consumidor e as intenções de compra de produtos sustentáveis. Esses construtos incluem posicionamento estratégico da marca, qualidade percebida do produto, inovação no varejo e autenticidade da marca.

Os resultados dos modelos de equações estruturais demonstram que todos os caminhos são estatisticamente significativos com forte poder de efeito. O posicionamento estratégico emergiu como o maior contribuinte para o efeito na disposição de comprar. Como hipotetizado, a equidade da marca e a sustentabilidade ambiental têm efeitos indiretos na intenção de compra comportamental. Os consumidores têm um alto nível de θ, enquanto os itens têm valores altos, corroborando a aceitação dos produtos, exceto para a diferenciação da marca.

O estudo também descobriu que o uso do discurso de inovação na estratégia de marca por si só pode não ter efeito na atração de consumidores. Para que as marcas analisadas alcancem ligações mais fortes na memória dos consumidores, é necessário ter uma adequação nominal para nomes que apresentem uma baixa frequência, uma maior persuasão em relação aos clientes sobre o discurso contido nas impressões de embalagens e outros materiais de mídia, e o desenvolvimento de embalagens sustentáveis.

As implicações teóricas dessas descobertas são substanciais. Elas se baseiam no conhecimento existente sobre equidade da marca e sustentabilidade ambiental, oferecendo uma compreensão diferenciada de como esses construtos interagem para influenciar o comportamento do consumidor. Esses dados podem ser aproveitados por acadêmicos e pesquisadores nas áreas de marketing e sustentabilidade para explorar ainda mais a dinâmica das relações consumidor-marca no contexto da sustentabilidade.
As implicações gerenciais deste estudo são igualmente importantes. Empresas, especialmente aquelas que atuam nos setores de cosméticos, turismo, vestuário, acessórios e alimentos, podem usar essas informações para posicionar estrategicamente suas marcas e produtos para atrair consumidores ambientalmente conscientes. Ao focar na qualidade do produto, inovação no varejo e autenticidade da marca, as empresas podem aumentar a equidade de sua marca e impulsionar a lealdade do consumidor e as intenções de compra.

Essas implicações são particularmente cruciais para pequenas empresas florestais que buscam alavancar a sustentabilidade ambiental para construir sua equidade de marca. Ao compreender os construtos que influenciam o comportamento do consumidor, essas empresas podem desenvolver estratégias de marketing eficazes que ressoam com seu público-alvo.

Os beneficiários desta pesquisa são diversos. Acadêmicos e pesquisadores podem utilizar essas informações para aprofundar suas pesquisas no campo da equidade de marca e sustentabilidade. Empresários, especialmente aqueles em pequenas empresas florestais, podem obter percepções valiosas sobre como incorporar a sustentabilidade em suas estratégias de marketing e \textit{Branding}. Formuladores de políticas também podem se beneficiar dessas descobertas. As percepções fornecidos por esta pesquisa podem informar a formulação de políticas que incentivam práticas sustentáveis, promovendo assim um ambiente que apoie o crescimento e o desenvolvimento de marcas sustentáveis.
 
Os consumidores também podem se beneficiar. Ao obter uma compreensão mais profunda das práticas sustentáveis das marcas que escolhem apoiar, os consumidores podem tomar decisões de compra mais informadas que se alinham com seus valores e crenças. Isso pode melhorar a experiência geral da marca, promovendo um sentimento de lealdade e satisfação.
Embora este estudo forneça visões valiosos sobre a relação entre consumidores e marcas sustentáveis, especialmente no contexto da Floresta Amazônica, ele também abre caminhos para pesquisas futuras. O foco em setores específicos, como cosméticos, turismo, vestuário, acessórios e alimentos, e a concentração geográfica na Floresta Amazônica, embora forneçam profundidade, também apresentam uma oportunidade para estudos futuros ampliarem o escopo e explorarem outros setores e regiões. Isso aumentaria a generalizabilidade dos resultados e proporcionaria uma compreensão mais abrangente da dinâmica do \textit{Branding} sustentável.

O uso de modelos de equações estruturais neste estudo, embora eficaz, pressupõe relações lineares e pode não capturar totalmente a natureza intrincada do comportamento do consumidor e das estratégias de marca. Isso apresenta uma oportunidade empolgante para pesquisas futuras empregarem modelos mais complexos ou abordagens qualitativas que aprofundem as nuances do comportamento do consumidor. Além disso, a exploração de fatores além do comportamento racional do consumidor, como emoções, influências sociais e vieses inconscientes, poderia proporcionar uma compreensão mais rica dos impulsionadores da lealdade do consumidor e das intenções de compra de produtos sustentáveis. Essas possíveis melhorias destacam a natureza em constante evolução da pesquisa nessa área e a busca contínua por visões mais profundos e abrangentes.
